\documentclass[10pt,a4paper]{amsart}
\usepackage[margin=19mm]{geometry}
\usepackage{amsmath,amssymb,mathtools}
\usepackage[scr=rsfs,cal=euler]{mathalfa}
\usepackage{xfrac}
\usepackage[unicode,hidelinks,bookmarks=false]{hyperref}
\newcommand{\ie}{i.e.\ }
\newcommand{\cf}{cf.\ }
\newcommand{\resp}{respectively\ }
\newcommand{\Subsection}{\subsection}
\providecommand{\nobreakdash}{\nobreak\hbox{-}}
\setlength{\emergencystretch}{2em}
\multlinegap0pt
\usepackage{bm}
\usepackage{bbm}
\usepackage{dsfont}
\usepackage{xspace}
\usepackage{cancel}
\usepackage{mathtools}
\usepackage{amsthm}
\makeatletter
\@ifundefined{theorem}{\newtheorem{theorem}{Theorem}}{}
\@ifundefined{lemma}{\newtheorem{lemma}[theorem]{Lemma}}{}
\makeatother
\newcommand{\cI}{\mathcal{I}}
\newcommand{\cM}{\mathcal{M}}
\newcommand{\dd}{\,\mathrm{d}}
\title[Maximal operators given by Fourier multipliers with dilation]{Maximal operators given by Fourier multipliers with dilation of fractional dimensions}
\author{Jin Bong Lee, Jinsol Seo}
\date{Source: arXiv:2405.16855v2 (CC BY 4.0)}
\begin{document}
\maketitle

\noindent\emph{Source and licence.} Maximal operators given by Fourier multipliers with dilation of fractional dimensions. Version arXiv:2405.16855v2 (CC BY 4.0), distributed under the Creative Commons Attribution 4.0 International licence, as recorded in the arXiv licence metadata for this exact version and shown on the abstract page https://arxiv.org/abs/2405.16855v2. Full rights notice: licenses/arxiv-2405.16855-CC-BY-4.0.txt.\par\smallskip

\noindent\textit{Editorial scope.} Counted unit: the Lemma whose source label is \texttt{231109742} in the article (source0.tex lines 615--625, source label \texttt{231109742}), delivered together with the article's own complete original proof of it (source0.tex lines 628--672, one \texttt{proof} environment of 45 lines). No mathematics is written, repaired or re-derived: statement and proof are the article's own text line for line. Required context retained uncounted: none. Every definition, display and label of those lines is retained unchanged. Omitted from this package and not redistributed: 1--614, 626--627, 673--1157. Nothing inside a retained excerpt is omitted or altered: every retained body is delivered exactly as the article's lines are written, so no omission receipt of any kind accompanies this package. Citations: the retained excerpts carry no citation command at all, so no bibliography entry of the article is transcribed and no reference is redistributed. Reference adaptation: none was needed and none was made; every reference inside the delivered unit resolves inside it, so no unresolved reference or citation is produced. Printed slips kept literal: no printed word, symbol, hypothesis or number of the counted statement or of its proof was altered. Layout and class: the article's own class is \texttt{amsart} with packages including amsmath, amssymb, mathrsfs, amsfonts, dsfont, amsthm, relsize, enumerate, verbatim, amscd, graphicx, appendix, tikz, hyperref; the wrapper is the lane's pinned \texttt{amsart} editorial wrapper at 10pt on A4, which restates the theorem-like environments the retained text needs and adds the article's own macro definitions that the retained text needs (\texttt{\textbackslash cI}, \texttt{\textbackslash cM}, \texttt{\textbackslash dd}), with extra wrapper packages (bm, bbm, dsfont, xspace, cancel, mathtools). The delivered numbering is the unit's own. Assets: no retained span contains a figure environment, a graphics-inclusion command, a PDF-page inclusion, a table or a TikZ picture; no image file is copied into this package, the package declares no asset dependency and no image file is redistributed with it. Licence: version arXiv:2405.16855v2 of this article is distributed under the Creative Commons Attribution 4.0 International licence, as recorded in the arXiv licence metadata for this exact version and shown on the abstract page https://arxiv.org/abs/2405.16855v2. A full rights notice is delivered at licenses/arxiv-2405.16855-CC-BY-4.0.txt. Characters: every retained excerpt is pure ASCII, so the delivered engine, XeLaTeX with its default Latin Modern text font, reports no missing character.
\begin{lemma}\label{231109742}
    Let $E$ be a non-empty set in $[1,2]$.
    For any $a\in(0,1)$, 
    \begin{align}\label{240313416}
        \sup_{0<\delta\leq 1}\delta^a N(E,\delta)\lesssim_a \int_1^2d(t,E)^{-1+a}\dd t\lesssim_a 1+\int_0^1 \lambda^{a}N(E,\lambda)\frac{\dd \lambda}{\lambda}\,.
    \end{align}
    In particular, the Minkowski dimension of $E$ equals
    \begin{align}\label{240409516}
        \inf\Big\{b>0\,:\,\int_1^2d(t,E)^{-1+b}\dd t<\infty\Big\}\,.
    \end{align}
\end{lemma}

\begin{proof}
Observe that
\begin{align}\label{231109722}
    \int_1^2d(t,E)^{-1+a}\dd t&=(1-a) \int_0^\infty \lambda^{-2+a}\big|\{t\in[1,2]\,:\,d(t,E)<\lambda\}\big|\dd \lambda\,.
\end{align}
For a fixed $\delta\in(0,1]$, we denote 
$$
\cI_\delta:=\big\{[k\delta,(k+1)\delta]\,:\,[k\delta,(k+1)\delta]\cap E\neq \emptyset\big\}\,.
$$
One can observe that if $I\in\cI_\delta$, then $d(\cdot,E)\leq \delta$ on $I$.
This implies that
\begin{align}\label{240422418}
&\delta N(E,\delta)=\bigg|\bigcup_{I\in \cI_\delta}I\bigg| \nonumber\\
\leq \,\,&\big|\{t>0\,:\,d(t,E)\leq \delta\}\big|\\
\leq \,\,&
(1-a)\delta^{1-a}\int_\delta^\infty \lambda^{-2+a}\big|\{t\in [1,2]\,:\,d(t,E)<\lambda\}\big|\dd \lambda+2\delta\nonumber\\
\leq \,\,&3(1-a)\delta^{1-a}\int_0^\infty \lambda^{-2+a}\big|\{t\in [1,2]\,:\,d(t,E)<\lambda\}\big|\dd \lambda\,.\nonumber
\end{align}
Here, the last inequality follows from that because 
$$
|\{t\in [1,2]\,:\,d(t,E)<\lambda\}|=1\quad \text{for all}\quad \lambda\geq 1\,,
$$
we obtain that for any $0<\delta\leq 1$,
$$
\delta\leq (1-a)\delta^{1-a} \int_1^\infty \lambda^{-2+a}|\{t\in [1,2]\,:\,d(t,E)<\lambda\}|\dd \lambda\,.
$$
Combining \eqref{240422418} with \eqref{231109722}, the first inequality in \eqref{240313416} is proved.

To prove the second inequality in \eqref{240313416}, 
observe that for a fixed $\lambda>0$,
$$
\big\{t\in[1,2]\,:\,d(t,E)<\lambda\big\}\subset \bigcup_{I\in\cI_\lambda}\widetilde{I}\,,
$$
where $\widetilde{I}$ is an interval with the same center as $I$ and $\ell(\widetilde{I})=3\ell(I)$.
Due to the definition of $N(E,\delta)$, we have
$$
\big|\big\{t\in[1,2]\,:\,d(t,E)<\lambda\big\}\big|\leq \bigg|\bigcup_{I\in \cI_\lambda}\widetilde{I}\,\bigg|\leq 3\lambda N(E,\lambda)\,.
$$
By combining this with \eqref{231109722}, we obtain that
\begin{align*}
&\int_0^\infty \lambda^{-2+a}|\{t\in[1,2]\,:\,d(t,E)<\lambda\}|\dd \lambda\leq \frac{1}{1-a}+3\int_0^1\lambda^{-1+a}N(E,\lambda)\dd \lambda\,.
\end{align*}
Therefore the proof of \eqref{240313416} is completed, and
it directly follows from \eqref{240313416} that $\dim_\cM E$ equals \eqref{240409516}.
\end{proof}

\end{document}
